% =====================================================================
%  Harald Høffding, Demokratisk Radikalisme. En Indsigelse
%  In: Tilskueren. Maanedsskrift for Literatur, Samfundsspørgsmaal og almenfattelige videnskabelige Skildringer,
%  udgivet af N. Neergaard.  Sjette Aargang.  Kjøbenhavn: P. G. Philipsens Forlag; Thieles Bogtrykkeri.  November--December 1889, pp. 849-872.
%  Antiqua.
%
%  TRANSCRIPTION -- GENERATED; DO NOT EDIT BY HAND.  Høffding's four pieces in Tilskueren 1886-1890 (Om Religionsfrihed;
%  the three contributions to the Nietzsche debate with Georg Brandes) are built by ONE pipeline
%  (ebook-method/tilskueren-1886-90/) into a combined working file, which split.py cuts into four transcriptions.
%  Rebuild from the repo root:   sh ebook-method/tilskueren-1886-90/rebuild.sh
%
%  STATUS: COMPLETE (2026-10-15): the volume's title page and the whole text, pp. 849-872, with its notes.
%
%  EDITION.  The only printing in Tilskueren.  Text of record: Google's scan of the Harvard copy, Internet Archive
%  tilskuerenmaaned1889nngoog, 600 ppi grey JP2 leaves with the Archive's ABBYY OCR; extract ~/bibliotek/Høffding, Harald/1889-tilskueren-6-pp849-872.pdf (SCANS.tsv, sha256),
%  made by ia_extract.py (title leaf, Indhold leaves, the article).  Page range confirmed from the volume's Indhold and the
%  signature.  Høffding's reply to Brandes's »Aristokratisk Radikalisme« (Tilskueren, August 1889).  Translated in full by William Banks, The Great Debate (University of Wisconsin Press, 2024); not translated here.  The catalog's earlier subtitle „En Indvending“ was wrong: the head, the Indhold and Runeberg's index read „En Indsigelse“.
%
%  ACCURACY (2026-10-15; one collation for all four texts, 76 pp., ~29,000 words).
%  (a) The DRAFT is the Archive's ABBYY OCR (the layer of the working PDF), turned into LaTeX by draft.py; it was collated
%      word by word with an independent tesseract reading (dan, tessdata_best) of the same images (collate/check.py,
%      witness 'ocr'): 650 disagreements.  check.py's own rules settled 44.
%  (b) A THIRD reading voted: Project Runeberg's OCR of the same Google images (never proofread there), by context
%      (rules.py umn-*, 372 decisions), guarded so that a vote never makes a word the corpus does not attest; other
%      rules (the layer's glyph faults against an attested word, accents tesseract adds, the ø rule, the house thin space)
%      settled 70.  The vote was AUDITED at the image: 40 rule decisions sampled blind and 30 more suspects read at zooms:
%      7 Runeberg-for-the-draft votes were wrong (ABBYY and Runeberg sharing a fault: bg/og, kim/kun, bar/har,
%      Anstrengelser/Anstrængelser, Vist/vist, én/en, h Vad/hvad), all mended; the ~30 letter-changing votes for the draft
%      not looked at are tesseract noise against a real word.
%  (c) The rest were read at the image BLIND (round b1: 303 crops on 31 sheets, six readers, 62/62 planted controls right):
%      203 open checks, 23 places where Runeberg disputed both OCRs, 54 places where the henrikpontoppidan.dk e-text
%      differs (a FINDING AID only: someone else's transcription, never a vote -- it normalises the print's spelling, e.g.
%      Nietzsche for the 1889 print's Nietsche), and the audit; ~150 decided in the main session from the readers' lines,
%      ~25 at 300-dpi zooms.  OPEN 0.  An unattested-word sweep of the collated text (353 words in no other transcription)
%      was read through; its doubtful words were looked at in zooms (one misprint found: forkellige).
%  (d) STRUCTURE read at the image (reader K): the title pages, the article heads, the section numerals (two the draft
%      had lost), the verse, the divider; every FOOTNOTE read whole (reader N: 11 notes; compared word by word with the
%      collated notes: 0 word differences, the reader right on the closing quotes).
%  (e) LETTERSPACING (the print's emphasis, set \emph{} here) is MEASURED at the image, not read: two blind line rounds
%      (l1, 256 lines; l2, the 134 doubtful ones again) agreed on only 40 of 134 lines, so spacemeasure.py decides by the
%      white between the letters of each word in tesseract's boxes (bimodal: 1-4 px roman, 6-8 px spaced; threshold 5 px;
%      isolated words must have every gap wide).  Checked at zooms on 16 words where the measure and the e-text's <em>
%      differ: the measure right on 11; three of its false positives (heavy isolated words: Ævne, frem, Side) led to the
%      isolation rule; three misses mended by hand after the page check (SPACEFIX in spaceapply.py).
%  (f) SAMPLED PAGE CHECK (pagecheck.py, seed 20261015; two blind readers): 6 whole pages, one or two from each piece
%      (Religionsfrihed p. 451; Demokratisk Radikalisme pp. 852, 868; Gensvar pp. 132, 140; Epilog p. 343), 2202 words:
%      1 WORD error (Talemaader for the print's Talemaaler: a correction wrongly taken from the e-text finding aid),
%      1 PUNCT (a speck read as '), both mended; LETTERSPACING 15 differences: at 300-dpi zooms 5 transcription errors
%      (tre, Kopier twice, Genmæle, fortsatte; mended), 4 where the transcription is right, 6 not decidable at 300 dpi
%      (the ends of runs in wide-set lines).  Word error rate after mending 0 in 2202 (95% upper bound ~1.4 per 1000);
%      letterspacing remains the least certain feature: ~2-3 per 1000 words.
%  PRINTED AS IS (misprints kept, marked at the line):
%    p. 864: „forkellige“ for forskellige; zoom z7
%  CONVENTIONS.  Page turns: % --- p. N --- and \opage{N} (the printed page); a word divided over a page turn is split with
%  %.  Footnotes as \footnote at their marks, numbered per page as printed (1)).  \emph{} = letterspacing; signatures in
%  small capitals as printed.  The running heads and the sheet signatures are not transcribed.  The ø rule (CLAUDE.md).
%  Quotation marks as printed: »...« and „...“ both occur, and the print sometimes doubles them (»»Zunächst).
% =====================================================================
\documentclass[12pt,a4paper,openany]{book}
\usepackage[utf8]{inputenc}
\usepackage[T1]{fontenc}
\usepackage{libertinus}
\usepackage{libertinust1math}
\usepackage[danish]{babel}
\usepackage[protrusion=true,expansion=true]{microtype}
\usepackage{setspace}
\usepackage[margin=1.3in]{geometry}
\usepackage{fancyhdr}
\usepackage{hyperref}
\hypersetup{pdftitle={Demokratisk Radikalisme}, pdfauthor={Harald Høffding}, hidelinks}
\pagestyle{fancy}
\fancyhf{}
\fancyhead[LE,RO]{\thepage}
\fancyhead[C]{\textit{H. Høffding: Demokratisk Radikalisme (1889)}}
\fancypagestyle{plain}{\fancyhf{}\fancyhead[LE,RO]{\thepage}}
\setstretch{1.3}
\usepackage{marginnote}
\usepackage{graphicx}
\DeclareUnicodeCharacter{0254}{\reflectbox{c}}   % ɔ (id est)
\renewcommand{\thefootnote}{\arabic{footnote})}   % the print numbers each page's notes 1), 2)
\newcommand{\opage}[1]{\marginnote{\footnotesize\textit{[#1]}}}
\newcommand{\sectionnum}[1]{\par\medskip\begin{center}#1\end{center}\smallskip}
\newcommand{\asterism}{\par\begin{center}*\quad*\\[-0.4em]*\end{center}}
\newcommand{\signature}[1]{\par\smallskip\hfill #1\hspace*{2em}\par}
\pdfstringdefDisableCommands{\renewcommand{\footnote}[2][]{}\renewcommand{\opage}[1]{}}

\begin{document}

\begin{titlepage}
\begin{center}
\vspace*{2em}
{\LARGE\bfseries \emph{TILSKUEREN.}}\\[1em]
{\large \emph{MAANEDSSKRIFT}}\\[0.5em]
{\small FOR}\\[0.5em]
LITERATUR, SAMFUNDSSPØRGSMAAL OG ALMENFATTELIGE\\
VIDENSKABELIGE SKILDRINGER.\\[1.5em]
{\small UDGIVET AF}\\[0.5em]
{\bfseries N. NEERGAARD.}\\[1.5em]
\emph{SJETTE AARGANG.}\\
\vfill
KJØBENHAVN 1889.\\[0.4em]
{\bfseries P. G. PHILIPSENS FORLAG.}\\[0.3em]
{\small THIELES BOGTRYKKERI.}
\end{center}
\end{titlepage}
\clearpage

% --- p. 849 ---
\setcounter{footnote}{0}%
\opage{849}
\begin{center}{\large\bfseries Demokratisk Radikalisme.}\\[0.4em]{\bfseries En Indsigelse.}\end{center}\medskip

\sectionnum{I.}

Naar jeg føler mig opfordret til at rette nogle Bemærkninger
mod den Artikel om »aristokratisk Radikalisme«, som Dr. Georg
Brandes offentliggjorde i »Tilskuerens« Augusthæfte, saa vil jeg
ikke, at det skal ske, uden at jeg først og fremmest og af oprigtigt
Sind udtaler, hvor meget jeg anerkender, hvad vor Literatur og
vort Aandsliv skylder denne Mands glimrende Begavelse og store
Ævne til at vække og begejstre. Jeg mener ikke, at hans Betydning
kan udtrykkes derved, at han har sat Problemer under
Debat; ti han har egentlig aldrig selv villet være Tænker; ---
hvad der har været hans Fortrin, er hans store Ævne til med
Energi og Varme, eller for at bruge et Udtryk, han selv i den
seneste Tid synes at ynde, med Voldsomhed at henpege paa saadanne
aandelige og literære Fænomener, som han anser det for
værdt at sysle med og lade sig paavirke af. Selv de, der ikke
paa deres eget Omraade have lært noget af ham, og som i mangt
og meget maaske endog have følt sig i Afstand fra ham, kunde
dog ikke andet end følge hans Virksomhed med Sympati og være
ham taknemlig for det Liv og den Bevægelse, han vakte. Jeg
ser her endda bort fra den vidunderlige Ævne, han besidder til
stemningsfuld Karakteristik, og som gør hans literaturhistoriske
Værker saa tiltrækkende.

Men den nævnte Artikel har baade ved hvad den angriber
og ved hvad den opstiller, grebet saa stærkt ind i hvad der for
mig staar som sandt og rigtigt, og den hele Tankegang, den hviler
paa, har i den Grad forarget mig, at jeg maa gøre, hvad jeg kan
% --- p. 850 ---
\setcounter{footnote}{0}%
\opage{850}for at bekæmpe den. Det forekommer mig tillige, at den bærer
Vidnesbyrd om det mislige i at indskrænke sig til en blot vækkende
eller æggende Virksomhed, i Stedet for at gaa selve de Problemer,
man rører ved, nøjere ind paa Livet. Jeg finder, at den Styrke
og Voldsomhed, hvormed Artiklen er skreven, staar i et underligt
Misforhold til Klarheden og Grundigheden i Tankegangen. Det
er ganske vist forbundet med nogen Vanskelighed at vide, hvem
man egentlig skal adressere denne Bebrejdelse til. Ti det er
ingenlunde altid let at afgøre, hvem de enkelte Udtalelser skrive
sig fra, Friederich Nietsche eller Georg Brandes. I sin Beundring
for den tyske Digter eller Tænker (hvad man nu vil kalde ham)
gør Artiklens Forfatter hans Tanker til sine og udtaler dem ofte
i sit eget Navn og i sit eget Sprog. Men det er ikke altid let at
vide, hvem man har for sig. En kritisk Undersøgelse af Forholdet
mellem dem vil jeg ikke indlade mig paa. Jeg anser mig berettiget
til at drøfte den »aristokratiske Radikalisme«, saaledes som den
her er gjort gældende, hvad enten Nietsche eller Brandes er Ophavsmand
til de enkelte Sætninger. Kun hvad et enkelt Punkt
angaar, vil jeg til Slutning gaa nærmere ind paa Nietsches Anskuelser
i og for sig.

Hvad jeg tror, Dr. Georg Brandes har forset sig paa, er
den Betydning, som Kraft og Hensynsløshed kan have, afset fra
Værdien af det, i hvis Tjæneste de staa. Der er noget tiltalende
og vækkende ved enhver kraftig Livsytring, ikke mindst paa det
aandelige Omraade. Den kan virke rensende og forfriskende, og
sætte andre Kræfter af større Værdi i Bevægelse. Men én Ting
maa dog aldrig glemmes: at selve Styrken ikke er noget Bevis
for Værdien eller Gyldigheden. Alligevel er det paa en Maade
dette, Nietsche hævder. Ti til de mange Ting, han polemiserer
imod, hører ogsaa Videnskaben, ja selve den videnskabelige Aand,
den Grundstemning og den Grundforpligtelse, hvorpaa al videnskabelig
Forsken hviler. Efter at have bekæmpet »det asketiske
Ideal« under forskellige Former, genfinder Nietsche det nemlig
ogsaa hos Videnskabsmændene. Disse have maaske aflagt mange
andre Fordomme, indrømmer han; men de tro endnu paa Sandheden,
og ere derfor ikke frie Aander. De hylde endnu ikke den
Maxime: »intet er sandt, alt er tilladt!« »Jeg kender«, fortsætter
Nietsche, »alt dette maaske altfor godt paa nært Hold: hin agtværdige
Filosofafholdenhed, til hvilken en saadan Tro forpligter,
hin Tankens Stoicisme, der tilsidst forbyder sig selv lige saa
% --- p. 851 ---
\setcounter{footnote}{0}%
\opage{851}strængt at sige Nej som at sige Ja\dots{}. Vi erkendende ere blevne
mistroiske mod alle Slags troende; en Mistro har efterhaanden
vænnet os til at slutte omvendt af, hvad man en Gang sluttede:
nemlig overalt, hvor en Tros Styrkegrad træder meget i Forgrunden,
at slutte til en vis Svaghed i Bevisligheden, til selve det antagnes
Usandsynlighed\dots{}. Hvis nu intet var mere guddommeligt end
Vildfarelsen, Blindheden, Løgnen?« (Zur Genealogie des Moral,
p. 167—169). --- Hvad jeg her er betænkelig ved, er ikke, at det
Spørgsmaal opkastes, hvilken Værdi videnskabelig Erkendelse har
for det menneskelige Liv. Et saadant Spørgsmaal er berettiget,
og jeg har selv søgt at give Bidrag til dets Besvarelse (se Afsnittet
om intellektuel Kultur i min Etik). Men naar den »Tankens
Stoicisme«, der vejer sit Nej lige saa vel som sit Ja, forkastes,
saa bortdrages derved Grundlaget fra al Videnskab. Det maa
nemlig vel mærkes, at et Ja i én Retning ikke nødvendigvis medfører
et Nej i en vis bestemt anden Retning; det kommer an
paa, om de to Retninger udelukke eller ophæve hinanden, og dette
maa nøje overvejes. Det drejer sig da her ikke om at tillægge
Videnskaben mere eller mindre Værdi; det gælder, om den i det
hele skal være til. Aldrig kan en Overbevisnings Styrke være
tilstrækkelig Garanti for dens Sandhed. Netop den Styrke, hvormed
den ytrer sig, vil gøre det vanskeligt at prøve de forskellige
Muligheder paa alsidig Maade; man vil være tilbøjelig til at holde
sig til den ene Mulighed, man nu én Gang har faaet fat paa med
saa stor Kraft. Det er sundt og godt og uundværligt for Livet,
at der danner sig stærke Overbevisninger; men de maa finde sig
i, at Kritiken (d. v. s. Prøvelsen) følger dem for at undersøge
Gyldigheden og Begrundelsen. Der er her to Sider af det menneskelige
Liv, den kraftige Selvudfoldelse og den besindige Prøvelse,
som det gælder at bringe i indbyrdes Harmoni. Det er vanskeligt,
og ingen løser Opgaven uden at komme til at hælde mere til den
ene eller til den anden Side. Men Livet er i det hele en vanskelig
Ting og bliver ikke lettere, fordi man slaar en Streg over en
væsentlig Side ved den Opgave, det stiller.

Den Bebrejdelse, Nietsche retter mod den videnskabelige
Aandsretning, er nu i Virkeligheden ubegrundet. Enhver kraftig
Livsrørelse har sin Ret, og Videnskaben lader den roligt udfolde
sig. Den staar som en foreløbig Magt, og Videnskabens første
Forhold til den er iagttagende: Lad os se, hvad Livet her bringer
os! Alt levende har jo sin Ret! --- Fra sin egen Erfaring kender
% --- p. 852 ---
\setcounter{footnote}{0}%
\opage{852}Forskeren Betydningen af de foreløbige Indskydelser og Hypoteser;
det er dem, som føre til Opstillingen af de Ideer, der saa siden
kunne prøves i det enkelte ved de strænge logiske og eksperimentale
Metoder. Det er usandt, at »Tankens Stoicisme« er mistænksom
overfor den kraftige Livsbevægelse, den levende Tro, eller at
den behøver at være det. Kritiken vaagner først, naar det Liv,
hvormed noget staar for Følelsen og Fantasien, forveksles med
objektiv Gyldighed. Saa kunne Konflikterne komme, indtil den
rette Begrænsning er naaet. Det er det dogmatiske, ikke det
levende og det kraftige, der udæsker Kritiken.

Efter disse Forbemærkninger gaar jeg over til, hvad der
egentlig ligger mig paa Hjærte, at begrunde en Indsigelse mod
den aristokratiske Radikalisme. Jeg vil først drøfte den Sætning,
som er opstillet i den ovenanførte Artikel: at Menneskehedens
Opgave er at frembringe store Mænd. Dernæst vil jeg belyse de
Indvendinger, der i samme Artikel er fremførte mod Velfærdsmoralen.
Til Slutning vil jeg undersøge Nietsches Lære om
Moralens Oprindelse.

\sectionnum{II.}

Det er ikke ganske let at overbevise sig om, hvad den
aristokratiske Radikalisme egentlig hævder. Der gør sig i Artiklen
øjensynligt en vis Usikkerhed gældende, som først og fremmest
maa fremdrages.

Der citeres p. 574 følgende Udtalelse af Nietsche: »Menneskeheden
skal uafbrudt arbejde paa at frembringe enkelte store
Mennesker --- \emph{dette og intet andet} er dens Opgave«. Og paa
næste Side hedder det, at »Nietsches Grundforestilling er, at et
\emph{Folk vil sige den Omvej}, Naturen gaar \emph{for at} frembringe en
halv Snes store Mænd«, --- »at Menneskehedens Maal ikke er et
i det uendelige udskudt, men maa ligge i dens højeste Eksemplarer«
»Masserne« (hedder det p. 581, uden at man véd, om man skal
tilskrive N. eller B. denne Ytring; NB vilde maaske være et passende
Tegn), »de er altid \emph{kun} at betragte som et af tre, \emph{enten som}
Kopier af de store Personligheder, daarlige \emph{Kopier}, udviskede
Kopier i daarligt Materiale, \emph{eller} som \emph{Modstand} mod de store,
\emph{eller} endelig som \emph{Redskab} for de store«\footnote[1]{Udhævelserne ere af mig, baade her og i det følgende.}. (Man mærke sig

1) Udhævelserne ere af mig, baade her og i det følgende.

% --- p. 853 ---
\setcounter{footnote}{0}%
\opage{853}dette Trilemma: enten Kopier, eller Modstand, eller Redskab.
Intet fjærde?)

Disse Udtalelser ere tydelige nok. Men en anden Gruppe af
Sætninger peger hen til, at de store Mennesker, de højeste Eksemplarer,
 selv dog igen ere Redskaber for mere omfattende Formaal,
saa at der synes at være et Maal, der er højere end det højeste!
Og disse Sætninger findes Side om Side med de første, saa at
selv den Læser, der ikke straks er mistænksom overfor store Ord,
føler sig tvivlraadig. Det hedder saaledes p. 574: »Det højeste,
Mennesket kan naa, er et heroisk Liv, et, i hvilket der kæmpes
med de største Vanskeligheder \emph{for noget, som paa en eller
anden Maade kommer alle til Gode}«, og »de Mennesker;
der ere Historiens Formaal«, betegnes som de ædle, \emph{de i stor
Stil gode, de}, i hvem det godes Genius virker«. --- Herefter
ligger det \emph{sidste} Maal dog ikke i de store Mennesker, men i det,
som de virke for. Og naar de kaldes »gode i stor Stil«, hvad
betyder det saa andet end, at de have Ævne til at meddele til
andre, --- at de i særlig Grad formaa at sprede Glæde og Lykke
Sundhed og Kraft, Skønhed og Sandhed omkring sig? Hvad kan
man ellers forstaa ved det Udtryk »god i stor Stil«?

Mellem de to Grupper af Sætninger maa der træffes et
Valg. For dem, der ikke veje deres Nej saa meget som deres Ja,
og som derved let komme til at sige baade Nej og Ja til samme
Ting, kan det maaske være vanskeligt at vælge her. Men Ja og
Nej er ikke god Filosofi. \emph{Dersom} nu den sidst fremførte Gruppe
af Sætninger fastholdes, bortfalder i Grunden hele min Betænkelighed
og hele min Kritik. Det er en Sandhed, jeg efter min
Opfattelse fuldt ud kan gaa ind paa, at det fremfor alt gælder
om at udvikle Mennesker, som i stort Omfang og med stor Kraft
kunne arbejde for, hvad der kan komme mange, maaske alle til
Gode. Ud fra Velfærdsprincipet kan denne Sandhed meget let
begrundes. Og hvilken anden Maalestok har man for Storhed, i
hvert Tilfælde for den Storhed, der skal være ét med storstilet
Godhed, end netop den: om der er Vilje og Ævne til at arbejde
for saa stor Velfærd som muligt, for saa mange, som muligt?

Men vælger man at fastholde denne sidste Gruppe af Udtalelser,
kommer man i Strid med hele Artiklens Tankegang og
egentlige Tendens. Og det ikke blot fordi de store Mennesker,
der erklæredes at være det højeste Maal, dog viste sig tillige at
være Midler og Tjænere, men ogsaa fordi Afstanden mellem dem
% --- p. 854 ---
\setcounter{footnote}{0}%
\opage{854}og »Massen« da ikke bliver saa stor, som den siges at være.
Selv de, der ikke kunne være gode i stor Stil, kunne dog være
det i mindre Stil. Ævnen til at meddele, til at sprede Lysglimt,
til at mildne, lindre og kvæge, mangler næppe hos nogen, selv
den usleste, naar blot hans Øje og Følelse kan oplades. De store
Mennesker gaa desuden \emph{kun} altfor ofte hen over Massernes Hoveder.
Deres Værk kan først ad saa mange Omveje komme de mange til
Gode, at hvis der ikke var en Skala af Storhed og Godhed, lige
fra de højeste Eksemplarer til de ringeste, vilde der intet opnaas
selv ved det ypperste Stormenneskes Virken. Og hvor begynder
saa den Masse, der kun er \emph{enten} Kopi \emph{eller} Hindring \emph{eller}
Redskab? Som det ofte har vist sig, at hvad der først antoges
for en Taagemasse ude i Himmelrummet, i Virkeligheden var en
Mangfoldighed af klare Stjærner, saaledes gaar det ogsaa med
Menneskemassen. Men ligesom der til at opløse Stjærnetaagerne
skal en god astronomisk Kikkert, saaledes skal der til at opløse
Menneskemassen gode og klare og milde Øjne. Et Øje, der selv
er tilsløret af æstetisk og intellektuelt Hovmod, af sygelig Lede
eller af den »adelsvældige Eneraadighed«, Nietsche og hans Tilhængere
lovprise, vil ikke kunne foretage en saadan Opløsning.
Lykkes den, finder man i ethvert Individ, selv det ringeste, et
higende og stræbende Jeg; man føler Livets Puls banke ogsaa her.
Ligesom Filosofen Heraklit sagde til de fremmede, der vilde
trække sig tilbage, da de traf ham siddende ved Køkkenilden:
»Kom kun ind; ogsaa her er Guder!« saaledes vil der her være
Lejlighed til at udraabe: »Ogsaa her er Mennesker!« Dette er
jo ogsaa det Feltraab, under hvilket Historien er gaaet frem.
Maalestokken har været, hvor mange der kunde være med til at
sige: »Ogsaa vi ere Mennesker!« Og dette Raab kan og skal
ikke forstumme, hvor stor en Gæld de mange med begrænset
Ævne og snævre Kaar end staa i til de faa store, der, som man
har sagt, ikke blot se Solen staa op før andre, men endog faa
den til at staa tidligere op og sprede sine Straaler videre om, end
den ellers vilde gøre. Ved at fremhjælpe det menneskelige Liv i
saa mange som muligt vil man desuden paa bedste Maade arbejde
for at tilvejebringe Muligheden for Stormenneskers Opstaaen.
Ved nærmere Undersøgelse har baade i Historien og i Naturen
det store og det smaa altid vist sig at hænge nøje sammen. De
levende Væseners Kaar ere indfiltrede i hverandre. Des urimeligere
er den Moloksdyrkelse, den aristokratiske Radikalisme anbefaler.
% --- p. 855 ---
\setcounter{footnote}{0}%
\opage{855}Den Bonde, der lærer at læse, kan maaske selv blive en »Renan
eller Littré«, eller han kan stifte et Hjem, der afgiver gunstige
Betingelser for en saadans Opvækst. Hvad der gør et Menneske
stort, vil tilsidst grunde sig paa et oprindeligt, ved Arvelighedskombination
frembragt Anlæg, i Forening med de Livsforhold og
Betingelser, under hvilke Væksten og Udviklingen finder Sted.
Der finder her en naturlig Udvælgelse Sted: selv de gunstigste
ydre Betingelser ere ikke nok, naar Anlægget mangler, men paa
den anden Side behøver Anlægget gunstige Betingelser for ikke
at dø bort under Kampen for Tilværelsen. Den Højde, til hvilken
der naas, beror altid paa det Niveau, fra hvilket der begyndes.
Og jo mere den menneskelige Udvikling --- til Dels netop ved
Stormenneskenes Virken --- skrider fremad, des større Indflydelse
faar Slægten paa den enkelte, des mangesidigere bliver Afhængigheden,
og des større Betydning faar det fælles Arbejde, i hvilket
de mange nedlægge deres Vilje og deres Ævne, i Sammenligning
med hvad enkelte, de være nok saa store, formaa.

Dog ville Stormennesker aldrig kunne undværes. Der vil
stedse være Opgaver, som kun en stor enkelt, ikke mange smaa,
kunne løse. Kun i et enkelt Menneskes Aand kan den Tilspidselse,
den Koncentration, den inderlige Fordybelse blive mulig, uden
hvilke det store nye ikke undfanges. Først naar denne Koncentration
i en enkelt Bevidsthed som i et Brændpunkt har fundet
Sted, kommer Sekternes, Partiernes og Komiteernes Tid, og derefter
sker da Udstraalingen i de større Kredse. Det gælder ikke
blot for Videnskab og Kunst, men ogsaa for sociale og politiske
Opgavers Vedkommende. Religionernes Historie bærer særligt
Vidnesbyrd herom. Men ligesom Straalerne maa være der, før de
kunne samles, saaledes maa de have gunstige Betingelser for at
kunne udøve deres lysende og varmende Virksomhed, naar de
igen udgaa fra Brændpunktet. Der kommer intet Lys og ingen
Varme, dersom en klam, uigennemtrængelig Taage danner Omgivelsen.
En stadig Given og Modtagen finder her Sted, og Slagordene
om Stormenneskenes Enevælde sprænge den levende, historiske
Sammenhæng. Netop de virkeligt store have den bestemte
Følelse, at de have \emph{modtaget} meget, at de staa i en stor Gæld,
og at de intet havde udrettet, uden at gunstige Betingelser havde
virket med. Ingen har været mere opfyldt af Anerkendelse og
Pietet end Goethe, og Bismarck har udtrykkeligt afvist den Tanke,
at han skulde kunne lave Historie. Den Følelse, som disse
% --- p. 856 ---
\setcounter{footnote}{0}%
\opage{856}Stormennesker her udtale, stemmer ganske med hvad Udviklingslæren
fører til at antage. Det store vokser ud af det smaa, og
for at omsætte sin store Kraft til virkeligt Arbejde, behøver det
atter Modtagelighed og Ævne til Medvirken fra det smaas Side.\footnote[1]{I „Study of Sociology“ giver Herbert Spencer en Kritik af Overvurderingen af de store Mænd. Darwin sluttede sig ganske til ham, og skrev i Brev til ham: „Jeg har aldrig troet paa store Mænds herskende Indflydelse paa Verdens Fremgang“. (Life and Letters. III p. 165).}% the note runs over from p. 7125 to p. 7126 (joined by draft.py)

Det gaar desuden med de store, som med os alle: de have
deres Dyders Fejl, og Fejlene vokse endda ofte i højere Grad end
Dyderne. Det er nu én Gang saa, at Kraft og Flersidighed i
Regelen udelukke hinanden; Idealet er at forbinde begge, men i
Virkeligheden naas det \emph{kun} i ufuldkommen Grad. En kraftig
Ensidighed betegner Storheden hos de fleste, der fortjæne Navn
af store. Hvis nu Slægtens Liv i det hele ikke ogsaa skal blive
ensidigt, og derved i Længden forkvakles, maa Stormenneskenes
Virksomhed suppleres. Der maa være en Modvægt imod dem, for
at der ikke skal blive for mange Kopier af dem. Det er derfor
af Vigtighed, at man ikke abdicerer overfor dem, at der ingen
Kultus etableres overfor dem, men at den Anerkendelse, der bliver
dem til Del, udspringer af en kraftig Selvstændighed og Selvfølelse,
endog hos dem, der skylde dem mest. Men det udelukkes, naar
man sætter Maalet kun i dem, ikke i noget som helst af de
andre »Eksemplarer«. ---

Vi nødes da til at opgive den aristokratiske Radikalisme
og gaa over til \emph{den demokratiske Radikalisme}, der hævder,
at det sidste Maal maa være Udviklingen af Velfærd (i Betydning
af den til sund og kraftig Livsvirksomhed knyttede Lystfølelse)
hos alle. Saa snart blot nogle staa tomhændede tilbage, er
Maalet ikke naaet. Og selv om Flertallets Velfærd opnaas, er det
dog en Ufuldkommenhed, ofte en skærende Mislyd, der ikke maa
overhøres, naar det sker med Lidelse for et Mindretal.

Man behøver ikke at frygte, at Middelmaadighed eller Ensformighed
skal blive Følgen heraf. Det er netop en af Grundbetingelserne
for Slægtens Velfærd, at der findes en høj Grad af
individuelle Forskelligheder indenfor den. Det er endog en Betingelse
for dens blotte Bestaaen, som Darwin har paavist. Variering
er en Nødvendighed for, at en Art kan bestaa, naar den er naaet
til en vis Udvikling og Talrighed. Jo mere forskellige dens Individer
% --- p. 857 ---
\setcounter{footnote}{0}%
\opage{857}nemlig ere indbyrdes, des flere forskellige »Pladser i Naturens
Husholdning« kunne de udfylde. Denne fremskridende Individualisering,
der i det menneskelige Samfund blandt andet viser
sig i Arbejdets Deling, er nødvendig for, at Artens Formering
skal kunne overskride en vis, ikke høj Grad. Hvad der gælder
for den blotte Bestaaen, gælder ogsaa om Muligheden af at opnaa
en højere Udvikling af Artens Ævner og Kræfter. Det var ganske
i Udviklingslærens Aand, naar Stuart Mill (paa et Tidspunkt, da
denne Lære endnu ikke var fremtraadt) ud fra Velfærdsprincipet
fordrede den størst mulige Frihed for individuelle Tilbøjeligheders
og Anskuelsers Udvikling. Menneskene ere saa forskellige, at de,
selv om deres sidste Maal i Hovedsagen er det samme, dog ville
nærme sig det ad forskellige Veje, og Ejendommelighederne ville,
naar de faa Lov at udvikle sig, gøre Livet rigt og kraftigt. Schiller
har træffende sagt, at naar enhver stræber efter at naa Fuldkommenheden
paa sin Maade, kommer den ene ikke til at ligne
den anden, skøndt enhver kommer til at ligne det højeste Ideal:

Keiner sei gleich dem Andern, doch gleich sei Jeder dem Höchsten!

Wie das zu machen? --- Es sei Jeder vollendet in sich.
Det gælder om, at Harmonien mellem Kraft og Flersidighed tilvejebringes
hos saa mange enkelte som muligt. Derved vil Slægtens
Liv som Helhed blive saa rigt som muligt. Det er et Maal, som
til enhver Tid fremtræder under en ny, en forskellig Form, men
som ingen Tid kan undslaa sig for. Den, som bidrager mest til,
at vi nærme os til dette Maal, --- den, der som Forbillede, eller
ved sit Ord eller sin Handling yder mest Støtte til Opnaaelse af
Maalet, ære vi højest og kalde vi stor. I ham vil tillige Livet
røre sig kraftigst og fyldigst; han vil kunne have den største
Lykkefølelse, den, at raade over et stort Fond af Kraft, som anvendes
til et Formaal, der ikke blot omfatter hans egen Person,
men uendelig mange flere. Men skøndt vi ære ham højt og prise
ham lykkelig for, hvad han har været og virket, sige vi ikke, at
han er det højeste Maal, og hvis han virkeligt er stor, vil han
ogsaa vredes, hvis vi sige det. Ti han vil vide, at Maalet er langt
større, at det stedse stilles paany, og at Livets Kamp aldrig ender.

Et Aristokrati vil der altid være. Det er en Misforstaaelse
af Demokratiets Grundtanke (en Misforstaaelse, som ganske vist
ogsaa kan findes hos »Demokrater«), at vi nogensinde skulde
kunne undvære Mennesker, som er et Hoved højere end andre.
% --- p. 858 ---
\setcounter{footnote}{0}%
\opage{858}Der maa være Grundlæggere, Ledere og Ordførere. Der maa være
dem, der undfange de nye Tanker og give Stødet til de nye
Værker. Hvad Demokratiet vil, er netop, at det Aristokrati, som
ikke kan undværes, virkeligt bliver et Aandsaristokrati, --- at dets
aandelige Overlegenhed mere og mere kan komme til at træde i
Stedet for den paa ydre Midler grundede Forrang. Dette Maal
(som dog selv igen tillige er et Middel) nærmer det sig til ved at
skaffe gunstige Livsbetingelser for saa mange som muligt. Derved
forøges, som vi have set, ikke blot Muligheden for, at der kan
opstaa Stormennesker, men ogsaa for, at de kunne finde den
Anerkendelse, Tilslutning og Modvægt, uden hvilke deres Virken
ikke bliver frugtbar og gavnlig. Et virkeligt Aandsaristokrati vil
ogsaa selv have Følelsen af ikke at være det sidste Maal, men at
være et Redskab. Det højeste er at være den store Tjæner.
Det gamle Ord beholder sin Sandhed (det er jo ogsaa sagt af et
Stormenneske): »Den iblandt eder, som vil være stor, skal være
de andres Tjæner«. Æren er den, at være det store Middel, at
være meget for de mange --- hvad enten disse vide af det eller ikke.

Det er vistnok en rent æstetisk Tendens, der fører til Overvurderingen
af de store Mennesker. Man beundrer den store Kraftudfoldelse,
det vældige Stød, den Afstand, der synes at være
mellem dem og de fleste andre Mennesker, den Hensynsløshed,
med hvilken de i deres Sags Tjæneste gaa frem. Beundringen
fører let til, at disse Ting, især Hensynsløsheden, prises i og for
sig. Ved den æstetiske Vurdering tage vi den enkelte Person
eller Egenskab ud af dens Sammenhæng; vi se navnlig bort fra
dens praktiske Virkninger. Richard IIIs hensynsløse Grusomhed
betragte vi æstetisk blot som en Kraftudfoldelse, som Udslag af en
stor Lidenskab, der løber Linen ud. Gør man en saadan Betragtningsmaade
til den eneste, er det intet Under, at man kommer til den
aristokratiske Radikalisme, skøndt ogsaa det ringe og det begrænsede,
det enfoldige og det stille har sin æstetiske Værdi og sin Tiltrækning.
I en vis Forstand gør alligevel Æstetikeren de store Mennesker
(og de store Uhyrer) til Midler, nemlig til Midler for sin egen
æstetiske Nydelse. Det er jo ogsaa særligt af Hensyn til vor Skønliteratur,
at Artiklens Forfatter anbefaler Studiet af Nietsches
Skrifter saa indtrængende. Om Gavnligheden heraf tiltror jeg mig
ikke nogen Dom; Literaturhistorien frembyder mange Eksempler
paa, at Bekendtskabet med en Forfatter har medført Virkninger,
som netop danne en Kontrast til hans egen Retning. Jeg kan
% --- p. 859 ---
\setcounter{footnote}{0}%
\opage{859}dog ikke tro, at vor Literatur ikke skulde kunne faa det kraftige
Stød, den maaske trænger til, uden at der samtidig stødes an
mod simple Sandheder og mod Principer, man ikke kan modsige,
uden at modsige sig selv. Jeg tror, at baade Kunsten og Livet
udvikle sig sundest og klarest uden for mange Forvridninger og
Raffinementer. ---

At de store Mennesker ere store ved at være de store
Midler, ses ogsaa deraf, at det meget ofte var heldigst, om de
ikke havde været nødvendige. Dette gælder i høj Grad om den
af Nietsche beundrede Napoleon. Men det gælder ogsaa om den
af Georg Brandes med større Ret beundrede Cæsar. Om denne
hedder det: »Han [Cæsar], efter hvem al højeste Magt endnu
den Dag i Dag bærer Navn, kunde alt, vidste og kendte alt,
hvad en Hærfører og Hersker af højeste Rang maa kunne og
kende« (p. 596). Men den egentlige Motivering af Beundringen
maa søges i følgende Sætning: »For alle Fremskridt, der lode sig
udføre i hine Dage, var hans Liv en Borgen« (p. 597). --- Er
det dog ikke en sørgelig Tid, hvor en enkelt Mand, selv om det
er en Cæsar, er eneste Borgen for Fremskridtet? Det var en Tid,
hvor Menneskene kun længtes efter Ro, Fred og Trøst. Man
havde mistet Tilliden til egen Kraft, til egen Ævne til at føre Livet
paa en selvstændig Maade. Man saa opad, --- opad til de enkelte
store Magthavere, der benyttede den almindelige Slaphed til at
herske paa Jorden, eller opad til overnaturlige Magter, som forkyndtes
af Kristendommen og de andre Religioner, der paa denne
Tid trængte frem fra Orienten. At Cæsar kunde blive Fremskridtets
eneste Borgen, havde en lignende Aarsag, som at saa mange
opgav Troen paa det naturlige Menneskeliv og opsøgte en overnaturlig
»Borgen«. Tidsalderens Præg fremtræder i begge Henseender
som det samme. Men havde det ikke været lykkeligere,
om Livet havde udviklet sig saaledes, at Kraften og Modet vare
blevne holdt oppe, saa at man hverken i borgerlig eller i etiskreligiøs
Henseende havde opgivet sig selv? --- Jo mere det lykkes
at indpode Menneskene den Tro, at de allerfleste af dem intet
ere værd, --- at de kun ere »Kopier, Modstand eller Redskab«, ---
des mere vil ganske vist Trangen til en Cæsar vokse. Men det
vil da ogsaa gaa, som det gik Cæsar. Det ædle Ideal, som man
har ment stod for hans Tanke: at være \emph{Lederen} for de mange
\emph{frie Mænd}, kunde han ikke virkeliggøre. Mommsen siger om
ham: »Fra sin tidlige Ungdom var Cæsar Statsmand i Ordets
% --- p. 860 ---
\setcounter{footnote}{0}%
\opage{860}dybeste Betydning, og hans Maal det højeste, som det er givet
Mennesket at sætte sig: den politiske, militære, aandelige og
sædelige Genfødelse af den dybt sunkne romerske og den endnu
dybere sunkne helleniske Nation«. Mod dette Maal skal han have
styret baade som Demagog, og som Usurpator. Men uagtet han (efter
den anførte Historieskrivers Udsagn) »maaske er den eneste blandt
de vældige for Herren, der hverken i stort eller smaat nogensinde
har handlet efter Lune og Tilbøjelighed, men uden Undtagelse
efter sin Regentpligt« [altsaa som den store Tjæner; ellers
kunde han ikke have Pligter], saa naaede han dog ikke Maalet.
Hvad han grundlagde, var og blev, som Kejserhistorien viser, et
Tyranni, efter Omstændighederne velmenende eller grusomt. Ingen
Livsspirer udstrøedes, og de forhaandenværende kvaltes efterhaanden.
Derfor afslutter ogsaa Mommsen, der dog i Cæsar ser en »ubeskrivelig«
[!] Storhed, sin Vurdering af ham med Ord, der kunne
tjæne som Motto for den hele Betragtning, jeg har søgt at gøre
gældende: »Efter samme Naturlov, ifølge hvilken den simpleste
Organisme er uendeligt mere end den kunstigste Maskine, er ogsaa
enhver nok saa mangelfuld Forfatning, der yder \emph{Spillerum for
et Flertal af Borgeres frie Selvbestemmelse}, uendeligt
mere end den mest geniale og humane Absolutisme; ti hin er i
Stand til Udvikling, altsaa levende; denne er hvad den er, altsaa
død«. --- Man ser her Forskellen mellem en Æstetikers og en
Historikers Vurdering.

Dr. Brandes skelner mellem to Klasser af revolutionære
Aander, af hvilke den ene føler sig tiltrukken af Cæsar, den anden
af Brutus. Skulde der ikke kunne gives en tredje Klasse, der
hverken i Usurpatoren eller i den sidste Repræsentant for Roms
udartede Aristokrati finder den højeste Storhed, og som netop
finder hin Tid saa ulykkelig, fordi man kun havde Valget mellem
disse to? ---

Det virkeligt grundlæggende fandtes ved Oldtidens Slutning
snarest i det, som bevægede sig i den store Masse, hos de mange
smaa, de foragtede og de lidende. Her kom noget frem, som
viste sig at have mere Livskraft i sig end Cæsars Værk. Af
Askesens, Dogmatismens og Overtroens Svøb har der udviklet sig
paa den ene Side en Selvfølelse, en Følelse af Værd selv hos det
ringeste Menneske, en Følge af den Vægt, Kristendommen lagde
paa hvert enkelt Menneskes Sjæl og dens Frelse, paa den anden
Side en inderligere og mere omfattende Menneskelighedsfølelse end
% --- p. 861 ---
\setcounter{footnote}{0}%
\opage{861}man hidtil havde kendt. \emph{Det} blev de nye Kræfter, der kunde
føre Slægten videre. De blev ikke skabte af Kristendommen. De
ædleste Strømninger indenfor den græske Filosofi førte dem med
sig; men de styrkedes, inderliggjordes og udbredtes ved Kristendommen.
De bedste Mænd paa Kejsertronen i Rom vare i deres
bedste Foranstaltninger besjælede af denne Aand. Men først under
helt nye historiske Forhold kunde den i fuldere Maal komme til
Gennembrud. Det er denne Aand, ikke Cæsars Aand, som er det
sande Fremskridts Aand, omfattende alt, det store som det ringe.

\sectionnum{III.}

Der fremføres i Artiklen (p. 584 ff.) en Kritik af Velfærdsmoralen,
som indledes med den Bemærkning, at »den engelske
Velfærdsmoral ikke har slaaet an i Tyskland«. Dette er i en vis
Forstand sandt, skøndt der foreligger dygtige Arbejder (som de
nyligt af Gizycki og Friederich Paulsen udgivne), der gaa i denne
Retning. Men for den, der med Kritik og uden at holde sig til
mere eller mindre poetiske Aforismer vil sammenligne, hvad den
engelske Skole har ydet med den Maade, paa hvilken det er blevet
optaget i Tyskland, ville Tyskerne sikkert ikke komme til at staa
med Palmer i Hænderne. Jeg troede hidtil, at alle mulige Misforstaaelser
af Velfærdsmoralen vare forekomne hos tidligere tyske
Forfattere. Men af den foreliggende Artikel ser jeg, at der endnu
var nogle tilbage.

»Velfærdsmoralens Maal«, hedder det, »er at skaffe Menneskene
saa megen Lyst og saa liden Ulyst som muligt. Men hvad
nu, i Fald Lyst og Smerte ere saaledes sammenknyttede, at den,
der vil have saa megen Lyst som muligt, maa tage en tilsvarende
Sum af Kvaler med i Købet? Det hedder i Clärchens Sang:
\emph{Himmelhoch jauchzend, zum Tode betrübt. Hvem} véd,
om ikke det sidste er Betingelsen for det første.« --- Jeg maa tilstaa,
at jeg ikke forstaar den Logik, af hvilken denne Indvending
er fremgaaet. Dersom Ulyst skulde være en nødvendig Betingelse
for Lyst, maa man naturligvis finde sig deri. Velfærdsprincipet
siger jo blot: frembring saa megen Lyst og saa liden Ulyst \emph{som}
muligt! I hvilken Grad det er muligt at naa Maalet, og hvor
store Omkostningerne ville være, derom indeholder det i og for sig
intet. Principet omstyrtes ikke ved Vanskeligheden af at gennemføre
det. Det sejrer, som Erfaringen viser, stedse paany over disse
% --- p. 862 ---
\setcounter{footnote}{0}%
\opage{862}Vanskeligheder. Jeg vil her beraabe mig paa Clärchens egen
Autoritet, idet jeg vil fuldføre Citatet ved at fremdrage de sidste
Linjer, i hvilke hun uddrager Resultatet af hele sin Stemningsrække:
\emph{Glücklich allein ist die Seele, die liebt}! Det .er netop
Sagen: der gives en Bevægelse i Sindet, et Liv i Følelsen, som
er forbundet med en saa dyb og stor Tilfredsstillelse, at den
stærke Svingning mellen Lyst og Smerte ikke ophæver den totale
Lykkefølelse, men netop styrker den. Der virke her to psykologiske
Love sammen. Den ene er den, ifølge hvilken Smerten (eller
Ulysten, Sorgen), naar den ikke stiger over en vis Grad, kan
danne Baggrund for Lystfølelsen og derved forhøje denne. Allerede
heri ligger tilstrækkeligt Motiv til at vælge saadanne Tilstande i
Stedet for dem, der ikke naa saa højt i Intensitet, skøndt de paa
den anden Side ere mere usammensatte Lysttilstande. Den anden
Lov er den, at der i selve Sorgen kan være noget tiltrækkende,
sandsynligvis fordi det stærke Liv, den kraftige Bevægelse, den
Anspændelse af Ævnerne, som følger med den, i og for sig medfører
en Tilfredsstillelse. Al Kraftanvendelse og Anspændelse, som
ikke er uforholdsmæssig, er forbunden med Lystfølelse. Midt i
Sorgen kan da en Lystfølelse ved selve Sorgens Styrke og Voldsomhed
gøre sig gældende, hvilket viser sig i, at man ikke \emph{vil}
udrives af den. Den organiske Udladning, som er forbunden med
enhver stærk Stemning; er et vigtigt Element her, selv om den
ikke skulde indeholde den fulde Forklaring. --- Da Auguste Comte
havde mistet den Kvinde, som øvede en saa afgørende Indflydelse
paa hans Aandsretning i den sidste Periode af hans Liv, udtalte
han en Gang i en Henvendelse, han rettede til sit Erindringsbillede
af hende: »Kun dig skylder jeg, at jeg ikke kommer til at forlade
Livet uden paa værdig Maade at have lært den menneskelige
Naturs bedste Følelse at kende\dots{}. Midt under de tungeste
Smerter, Følelsen kan medføre, har jeg aldrig ophørt at føle, at
\emph{det væsentlige for Lykken er at have Hjærtet opfyldt
paa værdig Maade, selv om det er af Smerte, ja selv af
den bitreste Smerte}«. Auguste Comte og Clärchen ere altsaa
ganske enige, og Velfærdsmoralen er enig med dem. Med den
store Livsbevægelse og Livsfylde følge ogsaa de store Omkostninger,
de store Kontraster og de store Svingninger. Det gælder da at
have Mod til at vove og vælge. Hvem har sagt, at det højeste
skulde kunne faas for en billig Penge? ---

% --- p. 863 ---
\setcounter{footnote}{0}%
\opage{863}En anden Indvending lyder saaledes: »Hvorfor Lykke for
det største Antal? Man kunde tænke sig: for de bedste, de
ædleste, de mest geniale, og man maa have Lov til at spørge,
om jævn Velstand og jævn Velvære er at foretrække for den
Ulighed i Kaar, hvis Brod tvinger Kulturen til stadig at stige«. ---
Ved Besvarelsen af denne Indvending kan jeg vise tilbage til de
ovenfor anstillede Betragtninger om Stormenneskene. Det er klart,
at der særligt bør sørges for dem, som ere velfærdsbringende i
største Omfang og i højeste Grad. Alt hvad der kan gøres for
dem, bringer tusendfold Frugt. De bedste, de store Tjænere, maa
der netop ifølge Velfærdsprincipet tages et særligt Hensyn til. Det
er i sin Orden, at i et Slag den menige Soldat ofrer sig eller ofres
for den store Feltherre, paa hvem Tusenders eller et helt Folks
Frelse beror. Velfærdsprincipet angiver kun, hvilken Maalestok
der skal anvendes ved Vurderingen af de enkelte Handlinger eller
Foranstaltninger; i de enkelte Tilfælde vil Fordelingen af Lyst og
Smerte blive højst forskellig, alt efter Forholdene. Paa lignende
Maade indeholder Aarsagsprincipet ikke noget \emph{om, hvilke} Aarsager
der i de enkelte Tilfælde virke, og hvorledes de kombineres for
at frembringe Virkningen. Vore almindelige Principer føre først
til bestemte Resultater ved at bringes i Forbindelse med bestemte,
givne Tilfælde.

Hvad Ulighed i Kaar angaar, saa kan den ganske vist være
en Brod, der driver fremad. Og der er jo heller ingen Udsigt til,
at denne Brod nogensinde falder bort. Foreløbigt har Velfærdsmoralen
ingen Grund til at tage denne Udsigt med i Betragtning.
Der er saa megen afgjort Uretfærdighed og Ubillighed at afhjælpe
indenfor den nuværende Samfundsorden, at det kan være tidsnok
at tales ved om Undværligheden af Ulighed, naar vi ere komne ud
over de værste Misligheder. \emph{Dersom} hin Brod ikke kan undværes,
saa vil jo i øvrigt netop Velfærdsprincipet fordre, at den bibeholdes.
Hvorledes kan der egentlig heri ligge nogen Indvending?

Historisk er Velfærdsprincipet særligt blevet proklameret i
Modsætning til de begunstigede Klassers Magthaveri. Om Jeremy
Bentham, den Mand, der mest af alle har draget praktiske Konsekvenser
af dette Princip, skrev hans Ven, James Mill: »Benthams
Meninger voksede ud fra én Rod, nemlig at Menneskeslægtens Vel
er det forpligtende Princip. Han brugte hele sit Liv til at anvende
dette Princip paa de vigtigste menneskelige Interesser: Lovgivning,
Statsforfatning, Kirkeforfatning, Opdragelsen og praktisk Moral.
% --- p. 864 ---
\setcounter{footnote}{0}%
\opage{864}Paa alle disse Omraader fandt han, at de manges Interesser stadigt
vare blevne ofrede for de faas. Med andre Ord: i Stedet for Dyd
havde Last været den herskende Kraft ved de menneskelige Anliggenders
Ledelse. At drage Sløret fra dette Ubillighedens Mysterium
og at vise de mange, hvorledes de vare blevne behandlede, saaledes
som Bentham uden Skaansel gjorde det, maatte saare Følelserne
hos de smaa Kliker, der holdt de nationale Interessers forkellige % PRINTED AS IS (forkellige: for forskellige; zoom z7)
Grene i deres Haand, og som med Raseri betragtede den Mand,
der beviste Nødvendigheden af, at den store Mængdes Interesser
beskyttedes overfor dem«.

Velfærdsprincipet er et demokratisk Princip, for saa vidt
det forbyder at glemme de manges Lidelser over de faas Nydelser.
Men det er ikke fjendtligt mod de faa store, hvis Væsen og Virken
føre Liv, Udvikling og Velfærd med sig for store Kredse.

Selv om Ulighed i ydre Kaar stedse i en eller anden Grad
og i en eller anden Form vil findes, saa længe Menneskelivet bestaar,
saa er den dog ingenlunde det eneste Motiv til Fremgang
i Kultur. Trangen til at bruge sine Kræfter, til at udvikle sine
Ævner og faa Afløb for hvad man føler sig at sidde inde med,
er ogsaa et vigtigt Motiv, og naar ydre Hindringer fjærnes, vil
det kunne spille en langt større Rolle end det hidtil har gjort.
Nu hindre ydre Forhold altfor ofte de enkelte i at blive opmærksomme
paa deres egne Muligheder, og selv hvor Driften er vakt,
kvæles den let af ydre ugunstige Omstændigheder. ---

En tredje Indvending lyder saaledes: »Det er Grundmodsigelsen
i denne Moral, at den fordrer og anbefaler en Forsagen
af Jeget, som finder Sted til Bedste for et andet Jeg«. --- Den,
som har opstillet denne Indvending, synes ikke at kende noget til
den engelske Skole. Ellers maatte han vide, at det netop er et
af Særkenderne for den Moral, som denne Skole har grundlagt,
at den ikke fordrer Forsagelse eller Opofrelse, uden hvor Tilfredsstillelsen
af den egne Trang vilde medføre saa meget des større
Ulæmpe for andre Individer. Den polemiserer udtrykkeligt med
Askesens og den umotiverede Resignations Moral, idet den fordrer
Begrundelse for enhver Ulyst, man paafører sig selv, lige saa vel
som for enhver Ulyst, man paafører andre. Og den hævder, at
i langt flere Tilfælde, end man hidtil har tænkt sig, er Individernes
Lykke solidarisk, naar man blot gør sig Umage for at finde de
fælles Betingelser. Bentham gaar endog saa vidt, at han undertiden
udtaler, at alles Vel bedst fremmes ved hver enkelts kloge
% --- p. 865 ---
\setcounter{footnote}{0}%
\opage{865}Egoisme. Og baade Bentham og Spencer have desuden vist,
hvorledes Sympati og Filantropi ofte skaffe sig Luft paa en højst
uheldig Maade, idet de forkæle og forkvakle deres Genstande og
gøre dem til hjærteløse Egoister. Lige saa lidt som Velfærdsprincipet
(hvad det skulde synes overflødigt at bemærke) fordrer,
at man uden Grund skal tilføje sig selv Lidelse, lige saa lidt fordrer
det, at man i hvert Øjeblik blot skal være optagen af at skaffe
andre Lyst. Jeg tillader mig her at gentage, hvad jeg et andet
Sted (i min Etik Kap. 7) har udtalt herom. »Dersom vi stedse
vilde rette vor Færd efter den Trang og de Fornødenheder, som
Menneskene i Øjeblikket have, saa vilde vi derved sænke Niveauet
for det menneskelige Liv og den menneskelige Kultur. Velfærdsprincipet
fordrer, at vi ikke sky Striden mod Fordomme og Træghed.
Det bedste, man kan gøre for andre, er ofte simpelt hen at lade
dem føle, at de endnu i deres Ønsker og Fornødenheder staa for
lavt og ikke stille tilstrækkeligt store Fordringer. Saaledes gaar
--- for at tage et enkelt Eksempel --- den store Kunstner ofte sin
ensomme Gang, uforstaaet eller endog miskendt af den store
Mængde. Og dog følger han Velfærdsprincipet, uden at han
maaske tænker derpaa, ved strængt at hævde Kunstens Fordringer
Han forøger Slægtens aandelige Kapital og skænker den en Kraft,
som i Fremtiden kan komme til at virke i store Kredse. Kun
en kortsynet Opfattelse og Anvendelse af Velfærdsprincipet holder
sig til den i Øjeblikket herskende Trang i Stedet for at se hen
til Slægtens varige Livsbetingelser og til de stadige Kilder for nyt
Liv og ny Virksomhed«. ---

Endelig hedder det i den i Artiklen fremførte Kritik af
Velfærdsmoralen: »Der gives ikke, eller kan ikke gives éns Handlinger;
enhver Handling er enestaaende i sin Art, saa at alle
Forskrifter kun gaa paa Handlingens grove Ydersider«. For saa
vidt der er noget sandt heri, og det er der ganske vist, fremdrages
her en Vanskelighed, der ikke særligt rammer Velfærdsmoralen,
men ethvert Forsøg paa at vurdere menneskelige Handlinger og
Forhold. Hvorledes bære Nietsche og hans Efterfølger sig ad med
at afgøre, om en Handling er »stor« eller ikke? Naar den er
enestaaende, kan den ikke bringes ind under noget som helst almindeligt
Begreb, heller ikke under Begrebet Storhed. Velfærdsmoralen
er dog her heldigere stillet, idet den tager sit Udgangspunkt
fra Handlingernes Virkninger, som dog altid til en vis Grad
maa kunne efterspores. I øvrigt behøver den ingenlunde at holde
% --- p. 866 ---
\setcounter{footnote}{0}%
\opage{866}sig til de enkelte Handlinger. Erfaringen viser, at visse Følelser,
visse Arter af Sindelag i særlig Grad afføde velfærdsbringende
Virkninger. Til disse indre Kilder for Handling søger den moralske
Vurdering at trænge ind; den ydre Handling har ofte en i etisk
Henseende ubestemt Karakter, og den afgørende Dom beror i
mange Tilfælde paa Forudsætninger, som kun den handlende
selv kan være vidende om. Men Velfærdsmoralen gør ogsaa i
sidste Instans den enkelte til sin egen højeste Autoritet og opponerer
(eller bør i hvert Tilfælde opponere) mod den gængse Forveksling
af moralske Domme med den offentlige Menings Domme.

\sectionnum{IV.}

Nietsche har i sit Skrift »Zur Genealogie der Moral«
(Leipzig 1887) udtalt en Opfattelse af de nu herskende moralske
Værdibestemmelsers Opstaaen, som det vil have sin Interesse at
gaa noget nærmere ind paa.

I Fortalen opstiller han følgende Problem: »Under hvilke
Betingelser opfandt Mennesket hine Vurderingsdomme \emph{godt} og
\emph{ondt}? og \emph{hvilken Værdi have de}? Hæmmede eller fremmede
de hidtil den menneskelige Fremgang (Gedeihen)? Ere de Tegn
paa Nødstand, paa Forarmelse, paa Udartelse af Livet? Eller ytrer
sig netop i dem Livets Fylde, Kraft og Vilje, dets Mod, dets
Tillid, dets Fremtid?« --- De to Spørgsmaal: om Moralens historiske
\emph{Oprindelse} og om den herskende Morals Værdi, gaa for ham
sammen til ét og besvares under ét, skøndt de i og for sig aldeles ikke
høre nødvendigt sammen. Hvilken saa de moralske Forestillingers
første Oprindelse maatte være, hvilke Motiver der end fra først
af have ført Menneskene til at kalde noget godt eller ondt, saa
afgøres derved ikke nødvendigvis den Værdi, de nu have, lige saa
lidt som min Personligheds Værdi beror paa, hvem der var min
Oldefader. Historisk udviklede Forhold og Institutioner kunne
senere hen faa en Betydning, der aldeles ikke kan udledes af den
Maade, paa hvilken de fra først af bleve til. Dette er noget, som
vi allerede finde indenfor den organiske Verden. Pingvinens Vinge
er historisk set et Flyveredskab, men dens Værdi beror paa, at
den kan bruges til Svømmeredskab. Om denne Forskydning af
Værdibestemmelsen, som i den nyere historiske og filosofiske Etik
spiller en stor Rolle, har Nietsche ingen Anelse. Og lige saa
uskyldig er han, naar han indlader sig paa en Omvurdering af
% --- p. 867 ---
\setcounter{footnote}{0}%
\opage{867}alle Vurderinger, eller rettere paa at berøve Værdibestemmelserne
deres Værdi (Entwerthung aller Werthe). Han ser ikke, at han
dertil behøver en Maalestok. Til en Omvurdering eller til Afførelse
af Værdier behøver man lige saa vel en Maalestok, som til den
oprindelige Vurdering. Men han gør ikke Rede for, hvilken han
anvender. Hans Paastande komme derved til at svæve i Luften ---
noget, de ogsaa ved andre af deres Egenskaber egne sig vel til.
Den Maade, hvorpaa han, som anført, har opstillet Problemet,
stemmer ganske med den Maade, paa hvilken en Tilhænger af
Velfærdsprincipet opstiller det, og paa hvilken Velfærdsmoralen
stedse har opstillet det (ti en saadan »Entwerthung« foregaar
uafbrudt og er blandt andet en væsentlig Del af den videnskabelige
Etiks Opgave). Og dog er Nietsche selv jo en Modstander af
Velfærdsprincipet! ---.

Saaledes forholder det sig med Præliminærerne. Lad os
nu se paa selve Teorien.

Efter Nietsches Antagelse hidrører Prædikatet »\emph{god}« ikke
fra dem, hvem Godhed bevistes. Det udsprang deraf, at de
lykkelige, de mægtige, de fornemme, de højt staaende følte sig
selv og deres Færd som noget af første Rang, i Modsætning [man
lægge vel Mærke til denne Modsætning!] til alt »lavt, lavsindet,
gement og pøbelagtigt«. For at udtrykke denne Modsætning
kaldte de sig selv »\emph{de gode}«. Det er altsaa »Afstandens
\emph{Pathos}«, som har ført til at danne og benævne Værdibestemmelser.
Et Fingerpeg i denne Retning finder Nietsche i de moralske Benævnelsers
Etymologi, der skal vise, at »fornem« og »ædel« er
det Grundbegreb, af hvilket »god« i Betydning af »sjæleligt fornem«
med Nødvendighed udvikler sig, ligesom »gemen«, »lav«, kommer
til at betyde »slet«. Man vilde, mener han, tidligere have fundet
denne Forklaring, dersom det ikke var paa Grund af de i den
moderne Verden herskende demokratiske Fordomme. Det var
Race- og Standsmodsætningen, som fra først af bestemte Vurderingen.

Men paa et vist Punkt i Udviklingen er der, ifølge Nietsche,
sket en Omvæltning i Anvendelsen af de moralske Benævnelser,
en hel Revolution i Vurderingsmaaden. Det var Jøderne, fra hvem
denne Bevægelse, denne \emph{Slaveopstand mod Herremoralen},
udgik og i Form af Kristendom udbredte sig over Verden. Hele
Verden tilbeder jo nu tre Jøder og en Jødinde! (Jesus, Petrus,
Paulus, Maria). De afmægtige, de overvundne, de lidende og
% --- p. 868 ---
\setcounter{footnote}{0}%
\opage{868}ulykkelige hævnede sig --- i deres Nag --- paa Sejerherrerne og
--- i deres Misundelse --- paa alle lykkelige ved at proklamere
de elendige som de i Sandhed gode, de lidende som de i Sandhed
fromme. Naget (das Ressentiment) blev saaledes skabende og
avlede Værdier --- ved at vende de oprindelige Værdier om.

Denne nu herskende Slavemoral udsætter den europæiske
Menneskehed for en stor Fare. Alt bliver under dens Indflydelse
tyndt, mat, godmodigt, tæmmet, svækket --- kort sagt: »forbedret«.
De lidende have sammensvoret sig mod de sunde og sejerrige og
bestræbe sig for at gøre dem lige saa syge i Sindet, som de selv
ere. Og en større Ulykke kunde dog ikke ske, end at de lykkelige,
de sjæleligt og legemligt mægtige begyndte at tvivle om deres
Ret til Lykken, og at man underkendte den fuldt klingende
Klokkes Forret for den mistonende og sprukne! Kun de lykkelige
ere jo Borgen for Fremtiden. Naar de skulle løse deres Opgaver,
saa maa de ikke vakle i deres Tro og deres Frejdighed, og de
kunne ikke sætte deres Tid og Kraft til paa at være Læger,
Trøstebringere og Frelsere for de syge. De to Klasser af Mennesker
maa holdes vel ude fra hinanden, og de syge bør helst have
Læger, der selv ere syge.

De lykkelige, de mægtige, der kunne leve Livet ud af egen
Kraft, uden Hensyn til andre, ere bedre og retfærdigere end de
lidende, som tage til \emph{Genmæle} og afvise deres Overgreb. De
\emph{aktive} ere aabne og \emph{frie, de reaktive}, hvis Færd ikke bestemmes
ud af dem selv alene, men ved Henblik (Nag eller
Misundelse) til andre, ere forud indtagne og ensidige i deres
Domme. ---

Et Par Udtalelser af Nietsche skulle endnu anføres til
Karakteristik af hans Opfattelse af den nuværende Tilstand. ---
»Hvo som ikke blot har Næse at lugte med, men ogsaa Øjne og
Øren, han mærker næsten overalt, hvor han nu til Dags kommer,
noget som Galehus- og Sygehusluft\dots{}. De \emph{sygelige} ere Menneskets
store Fare\dots{}. De fra først af forulykkede, faldne og nedbrudte,
disse, de \emph{svageste}, er det, der især underminere Livet blandt
Mennesker, forgifte vor Tillid til Livet og til Menneskene paa den
farligste Maade. --- Og derfor god Luft! god Luft! I hvert
Tilfælde bort fra alle Kulturens Gale- og Sygehuse! Og derfor
godt Selskab, \emph{vort} Selskab! Eller Ensomhed, om saa skal være!
Men i hvert Tilfælde bort fra den indre Fordærvelses og det
hemmelige Syge-Ormebids slemme Dunster! --- Og det, for at vi
% --- p. 869 ---
\setcounter{footnote}{0}%
\opage{869}dog, mine Venner, en Stund endnu kunne forsvare os mod de to
værste Syger, der kunne være os forbeholdt, --- mod den store
Væmmelse ved Menneskene, og mod den store Medlidenhed med
Menneskene!« (Zur Genealogie der Moral p. 131---135). ---

Saa vidt Nietsche. --- Lad mig begynde med at sige, at
hans almindelige Tanke om de lykkeliges Ret og om Vigtigheden
af, at de ikke tabe Troen paa dén, forekommer mig rigtig og
betydningsfuld. Jeg har ogsaa selv udtalt noget lignende, ganske
vist paa en mere stilfærdig Maade (Etik p. 50. 110. 138 f.). Lykken
er ikke saa almindelig i Verden, at der ikke skulde være al Grund
til at frede om den med en vis Ærefrygt, hvor den findes, hvis
den da ikke grunder sig paa Tankeløshed eller Hjærteløshed. Og
naar vi betænke, at den lykkelige ogsaa ofte raader over Muligheder
til videre Fremgang for store Kredse, Muligheder, de ulykkelige
oftest mangle, saa faar denne Betragtning en saa meget des større
Vægt. --- Men er der egentlig Grund til at udtale den paa en
saa larmende Maade? Ligner Europa virkelig et stort Hospital?
Skulde den Fare, Nietsche befrygter, virkelig være til Stede? Ere
vor Tids Mennesker virkeligt saa optagne af Lidelserne rundt om
dem, at de lade sig deres egen Lykke berøve? Saa vidt jeg kan
skønne, nærmer Europa sig til at blive en stor Kaserne, en stor
Fabrik og et stort Tivoli; men Hospitalerne hæmme endnu ikke
Livet eller Livsglæden; man affinder sig temmelig let med dem,
der bebo dem. Og man affinder sig ogsaa temmeligt let med
dem, der føle Savn og Lidelser udenfor Hospitalerne, og er snarraadig
til at finde paa farisæiske Udveje til »med god Samvittighed«
at vende dem Ryggen. (Se Kjelland's Novelle »En god Samvittighed«).
Medlidenheden plager os sandelig ikke, selv om der
gives megen tankeløs Filantropi og rastløs »Velgørenhed«, der
snarest er at betragte som en Slags aandelig og legemlig Motion
og ingen hastemte Toner behøver at fremkalde. Nietsche kæmper
mod et Fantom, som maaske til Dels er fremkaldt ved hans egen
Angest for den lurende Sygdom, men som i det mindste for en
Del skriver sig fra hans Studium af Schopenhauer. Han er aldeles
ikke, som Dr. G. Brandes siger, frigjort fra Schopenhauer; han
ser endnu paa alt med Schopenhauerske Briller. For Schopenhauer
var Væmmelse den herskende Følelse ved Verden og Menneskene,
og Medlidenhed den Følelse, der opstod som Modsætning og Lægedom.
Nietsche gaar nu ud fra, at Schopenhauer har sejret, og
at »Ekel und Mitleid« have delt Verden imellem sig. Han har
% --- p. 870 ---
\setcounter{footnote}{0}%
\opage{870}de Schopenhauerske Tanker, kun at han stiller noget omkring
med dem, og at han ikke blot føler Væmmelse ved Verden, men
ogsaa ved --- selve Medlidenheden!

Nietsche mangler ganske Blik for Storheden i den Anerkendelse
af de lidendes Menneskerettigheder, som Buddismen og
Kristendommen gennemførte, skøndt de tillige paa asketisk Vis
overdrev dem. Spirerne til denne Anerkendelse laa langt tilbage
i den indoeuropæiske Menneskeheds Historie. Det er ingen Slaveopstand,
der har ført til den. Allerede i den gammelariske Familje
gjaldt der en hellig Forskrift om at ære den fremmede, den- nødlidende
og den bønfaldende. Man har sikkert med Rette (saaledes
for ganske nyligt \emph{Leist} i hans Værk: Altarisches Jus Gentium.
Leipzig 1889) heri fundet Spiren til Humanitetsfølelsen. »Den
Sætning, at man for at undgaa Straffe fra Gæsterettens Guder
straks maa modtage Gæsten, Tiggeren og den bønfaldende med
Venlighed, har til sit Grundlag den Tanke, at man trods alle
Forskelligheder og Skæbnens Omskiftelser i Mennesket stedse skal
anerkende Mennesket, altsaa sin Lige, og at man derfor ogsaa
skal behandle det menneskeligt«. (Leist p. 379). Hos Grækerne
finde vi det samme, I Odysseen siges de fattige og bønfaldende
særligt at høre Zeus til. Altsaa ikke de lykkelige, men de ulykkelige
stilles i et særligt Forhold til Zeus, og det hos det Folk, som man
ellers plejer at opfatte som den lyse og harmoniske Livsglædes
Folk! Og den Sætning, at det er bedre at lide Uret end at gøre
Uret, en Sætning, som efter Nietsche maa høre ind under »Slavemoralen«,
er første Gang udtalt af Platon, den store græske Tænker
og den rige ateniensiske Adelsmand! --- Ikke en Slaveopstand,
men en Uddybelse og Udvidelse af Følelseslivet har ført til en
Ændring i Værdibestemmelserne, der i mange Punkter førte til en
afgørende Modsætning til den gamle militaristiske Moral. Den
hele fredeligere Udvikling af Livet medførte Erfaringer, man medens
det krigerske Liv herskede, ikke havde kunnet gøre. Og derved
blev der ogsaa vakt en Selvfølelse hos dem, der hidtil ikke antoges
at have noget Selv. ---

De »aktive«, siger Nietsche, ere bedre end de »reaktive«.
Ja, Undertrykkerne kunne naturligvis lettere lade være at tænke
paa de undertrykte, end disse paa hine. For saa vidt vil disses
Færd ikke blive bestemt saa meget ud fra dem selv som hines.
Magthaverne se vel de usle og svage færdes om sig, men behøve
ikke at betragte dem paa anden Vis end som ligegyldige Fænomener.
% --- p. 871 ---
\setcounter{footnote}{0}%
\opage{871}En særegen Animositet overfor dem er der ikke Grund til Men
den meste Grusomhed, der er øvet i Verden, er maaske netop
Ligegyldighedens Grusomhed. Naar de usle og afmægtige staa i
Vejen, naar de synes at ville brede sig, sparker man dem bort,
uden at tænke videre paa dem. Optagen af sin egen Glæde eller
Sorg, som vil have Luft, vil man ikke forulæmpes ved Synet af
disse Væsener. Jeg forklarer mig paa denne Maade, at Tor ved
Balders Baalfærd sparker Dværgen Litr, som aldeles intet havde
gjort, ind i Ilden. Her er ganske vist intet »Ressentiment«. Men
efter hvilken Maalestok kalder man slig Færd »bedre« end den,
de svage udvise for at hævde sig overfor de mægtiges Overgreb? ---
Desuden er ifølge Nietsches egen Antagelse de moralske Værdibestemmelser
hos de mægtige opstaaet ved en \emph{reaktiv} Følelse.
Eller hvilket andet Navn vil man give den Følelse, som han kalder
\emph{Afstandens Pathos}, og som skal have været den ledende ved
Moralens Oprindelse? Kan man føle sig i Afstand fra andre, uden at
denne Følelse bestemmes ved et Henblik til dem? Den oprindelige
moralske Følelse skal jo efter Nietsche være den ved \emph{Nedseen} paa
visse andre bestemte Selvfølelse. --- Atter et Ja og et Nej, der
mødes i hans Teori paa en skæbnesvanger Maade.

Men det er nu utvivlsomt aldeles urigtigt, at de moralske
Følelser og Forestillinger først og fremmest skulde være bestemte
ved en saadan Afstandsfølelse. Nietsche hænger sig her for meget
i visse Etymologier, der, selv om de ere rigtige, ikke bevise noget
afgørende. Langt større en Rolle end Forholdet til andre Racer
eller Stænder maa dog paa primitive Udviklingstrin Forholdet til
de andre Medlemmer af samme Klasse eller Familje have spillet.
Under den fælles Kamp for at bestaa fører dels Sympatien, dels
den gensidige Fordring om Hjælp og Bistand til en stadig Betragtning
og Vurdering af Handlingerne, de andres og de egne.
Det bliver Genstand for alles Opmærksomhed, i hvilken Grad de
fælles Interesser ske Fyldest. Det kommer an paa, hvem der er
den bedste Kriger og Jæger, den kløgtigste Medicinmand, den
snildeste Raadgiver. Her var Stof nok til Vurdering. Indenfor
enhver Gruppe af Mennesker er der Forskelle i denne Henseende;
alle ere ikke éns, yde ikke lige meget. Der var Pligter, der stadigt
trængte til at indskærpes selv indenfor de frie Krigeres Kreds;
der var naturlige Autoriteter (Forældre), eller der dannede sig
uvilkaarligt Autoriteter, som kunde indskærpe dem. Der dannede
sig snart en offentlig Mening, og der udtaltes Domme til Ære og
% --- p. 872 ---
\setcounter{footnote}{0}%
\opage{872}til Vanære. Hvis bonus, som Nietsche mener, egentligt betyder
»Kriger«, saa var »Kriger« et Æresnavn, som den fejge gik tabt
af, medens han til Gengæld fik et andet Navn (f. Eks. »Flygtningen«,
smlgn. Herodots Fortælling om Spartaneren Aristodemos, der paa
Grund af Sygdom vendte hjem og derved ikke kom til at deltage
i Kampen ved Thermopylæ). Selv en saa drabelig Helt som Sigurd
Fafnersbane faar (i »Sangen om Sigdrive« i den yngre Edda) Formaning
om at hævne sig paa den, der fornærmer ham, og ikke
vise sig fejg. --- Al moralsk Vurdering er reaktiv, for saa vidt
den forudsætter en \emph{Forskel} mellem et Ideal og en ringere Virkelighed.
Men denne Forskel eller Afstand maa komme frem allerede
indenfor samme Klasse eller Race; Forholdet til andre Klasser
eller Racer ligger her fjærnere. Den Afstandsopfattelse, Nietsche
appellerer til, er for reaktiv. ---

Dersom det ikke ved Artiklen om »aristokratisk Radikalisme«
var paanødt mig at betragte Nietsche som Filosof, vilde jeg ikke
have anstillet denne Kritik. Nietsche er ganske sikkert en sjældent
begavet Forfatter, men hans Begavelse er ikke filosofisk. Han
mangler den filosofiske Interesse og den filosofiske Resignation.
Havde han haft dem, havde han vel ikke egnet sig til at serveres
for Æstetikere. Han er Digter, Retor, hans Styrke er Stemningen
(særligt den antipatiske Stemning), som han formaar at udtale
med mærkelig Kraft. Men bag det hele mærker man, forekommer
det mig, det sygeligt forvredne. Vemodigt er det at tænke paa,
at han, som har formaaet at tage Ordet ofte med saa kraftig
Pathos (selv om det var »Afstandens« Pathos), at han nu selv er
en af de »sprukne Klokker«, hvis Mislyd han har raset saa
heftigt imod!
\signature{\textsc{Harald Høffding}.}

\end{document}
